\documentclass[12pt]{article}

% House style preamble
\input{.house-style/preamble.tex}

% Update PDF metadata
\hypersetup{
  pdftitle={Determinatives as nouns}
}

\title{Determinatives as nouns}
\author{Brett Reynolds \orcidlink{0000-0003-0073-7195}%
\thanks{Contact: \href{mailto:brett.reynolds@humber.ca}{brett.reynolds@humber.ca}}\\
Humber Polytechnic \& University of Toronto}
\date{\today}

\begin{document}
\maketitle

\begin{abstract}
English determinatives are a sub-category of the lexical category \term{noun}, alongside common nouns, proper nouns, and pronouns. The argument preserves \textit{CGEL}'s category/function distinction and the N-headed NP analysis. \textit{CGEL}'s existing treatment of independent genitives (\mention{mine}, \mention{Kim's}) shows that fusion of Det and Head functions, where it applies, is compatible with noun heads; the paper extends this conclusion to the determinative class. Partitive constructions (\mention{some of mine}, \mention{many of them}) establish that determinatives are head-like pivots in matrix NPs, and three diagnostics — paradigm productivity, structural saturation, internal modification — show that simple cases (\textit{Some left}) pattern with non-fused noun-headed NPs. The categorial claim doesn't depend on resolving the fusion question for every construction. Two simplifications follow: the DP/genitive-NP disjunction at the determiner slot collapses to NP, and the fusion-of-functions apparatus narrows to a smaller residue. The articles \mention{the} and \mention{a} are handled as defective members. The scope is synchronic English.
\end{abstract}

\section{Introduction}\label{sec:intro}

English has four sub-categories of the lexical category \term{noun}: common noun, proper noun, pronoun, and determinative. The first three are sub-categories of \term{noun} in \textit{CGEL} \citep{huddleston2002}; this paper argues that determinatives satisfy the same criterion for nounhood and should be added as the fourth.

The proposal extends an existing tradition. \textcite{spinillo2000determiners,spinillo2004reconceptualising} argue on English data that the determiner class isn't a valid primary lexical category. \textcite{vaneynde2003determiner} argues within HPSG that determinatives should be reclassified across the A and N categories on inflectional evidence from Italian and Dutch. Neither paper commits to the specific architecture defended here: all four of common, proper, pronoun, and determinative as coordinate sub-categories of a single supercategory \term{noun}.

The kind of taxonomic move has precedent in the same descriptive framework. \textcite{pullumwilson1977} dissolved the AUX category into V; \textit{CGEL} preserves this, treating modal and primary auxiliaries as a sub-category of verb rather than members of a separate functional category. The present proposal extends the same architectural move to determinatives within \term{noun}. The precedent isn't an argument from analogy but a demonstration that the framework has absorbed a closed functional class into an open lexical category before.

The paper keeps the category/function distinction and the N-headed NP analysis that \textcite{pullummiller2022nps} defend against DP. The difference from \textcite{abney1987} lies in the supercategory classification of determinatives. On the present account they're nouns with a distinctive sub-cluster profile, not heads of a separate functional projection.

The categorial claim doesn't depend on abolishing the fused Det-Head analysis. \textit{CGEL}'s existing treatment of independent genitives (\mention{mine}, \mention{Kim's}) already shows that fusion, where it applies, is compatible with noun heads. The paper extends this to the determinative class. Two stages of evidence support the extension. First, partitive constructions like \textit{some of mine} establish that determinatives are head-like pivots in matrix NPs, independently of whether the construction is fused. Second, three discriminators — paradigm productivity, structural saturation, internal modification — show that simple cases like \textit{Some left} pattern with non-fused noun-headed NPs and unlike adjectival Mod-Head fusion. \textit{CGEL} can keep fused Det-Head for partitive and similar cases without disturbing the categorial conclusion: the items doing the fusing are nouns either way.

The articles \mention{the} and \mention{a} are the worst case: they never head NPs in any structure. The paper handles them as defective members of a sub-category whose profile otherwise satisfies the criterion, on the same pattern as other lexically defective members within closed sub-categories (the suppletive \mention{none} for \mention{no} among determinatives, defective case forms among pronouns). Section~\ref{sec:articles} develops the account.

Two downstream simplifications follow. The DPs and genitive NPs that function as Det — currently a disjunction of phrase types — collapse to NP when determinatives become nouns. Other phrases that function as Det (PPs like \mention{about twenty}) are unaffected, but the primary noun-headed disjunction is simplified. If the simple cases are ordinary head-of-NP function rather than fusion, the fusion-of-functions apparatus developed in \textcite{Payne2007} drops out of most of its \textit{CGEL}-Ch.~5 applications.

§\ref{sec:bridge} shows that the determiner function is already routinely filled by NPs headed by nouns in \textit{CGEL}'s existing analysis (genitive pronouns, genitive common and proper NPs). §\ref{sec:argument} makes the central argument, first separating partitive headhood from the fusion question and then defending the simple-case reanalysis in §\ref{subsec:fused-head}; §\ref{subsec:det-uniform} develops the determiner-slot simplification. §\ref{sec:articles} handles \mention{the} and \mention{a}. §\ref{sec:nearby} distinguishes the present proposal from nearby work. §\ref{sec:subcluster} develops the sub-cluster differences internal to \term{noun}. The scope is synchronic English.

\section{Bridge from Noun to Determiner}\label{sec:bridge}

Before §\ref{sec:argument} makes its case, one fact about \textit{CGEL}'s existing taxonomy bridges the proposal. The determiner function is already routinely filled by NPs headed by nouns. \textit{CGEL} accepts this pattern in two recognized sub-cases.

First, NPs headed by genitive pronouns (\mention{my}, \mention{your}, \mention{his}, \mention{her}, \mention{our}, \mention{their}, \mention{its}) fill Det function:
\ea\label{ex:pronoun-det}
\ea \textit{my book}
\ex \textit{her office}
\ex \textit{their proposal}
\z
\z
Since pronouns are already a sub-category of \term{noun} in \textit{CGEL}, these NPs are ultimately noun-headed: the pronoun heads a Nom, and the Nom heads the NP.

Second, genitive NPs headed by common or proper nouns fill the same function:
\ea\label{ex:gen-NPs}
\ea \textit{Kim's book}
\ex \textit{the king's every whim}
\ex \textit{that man next door's car}
\z
\z
Each is an NP with its own internal modifiers and its own head N \citep[2]{pullummiller2022nps}.

Across these two sub-cases, the determiner function is already filled by NPs headed by members of three sub-categories of \term{noun} under \textit{CGEL}'s existing analysis: pronoun, common noun, and proper noun. The residue — items in Det function not captured by these two sub-cases — is the DPs headed by determinatives (\mention{the book}, \mention{this book}, \mention{some books}, \mention{every book}). The paper argues in §\ref{sec:argument} that determinatives are nouns too, which makes their \enquote{DPs} NPs. The determiner function then admits noun-headed NPs from all four sub-categories of \term{noun}, in addition to the non-noun-headed phrases \textit{CGEL} already accepts in Det function (PPs like \mention{about twenty}).

\section{The category-membership argument}\label{sec:argument}

The proposal needs a clear test for nounhood. The membership criterion is this: a lexical category C is a sub-category of \term{noun} iff items of C can head Nominals heading NPs that fill subject, object, or preposition-complement function, where Head is part of the category's profile rather than a function reached via fusion with another function. The criterion respects \textit{CGEL}'s phrasal architecture: a lexical N heads a Nom, the Nom heads an NP, and the NP fills subject (or some other) function. Subject isn't directly a function of N, and N doesn't directly head NP — the Nom layer is the intermediate constituent.

The criterion applies at the sub-category profile level, not at the level of individual lexemes. A sub-category qualifies if its items canonically head NPs as part of the sub-category's profile, and lexically defective members are admitted on paradigmatic grounds. Lexical defectiveness within a sub-category is familiar in the grammar: the auxiliary verb \mention{must} lacks a past-tense form (the suppletive \mention{had to} fills the gap), the determinative class includes suppletion (\mention{none} for \mention{no} in fused Det-Head uses; \textcite[22--23]{Payne2007}), and common nouns show mass-versus-count defectiveness. The same allowance covers \mention{the} and \mention{a} within the determinative sub-category, defective members of a sub-category whose profile satisfies the criterion. Section~\ref{sec:articles} develops that point. Sub-cluster differences in discourse profile, modifier selection, and inflection are developed in §\ref{sec:subcluster}.

\subsection{The fused-head reanalysis}\label{subsec:fused-head}

The central question is whether occurrences of determinative-headed phrases as Head in fused-head constructions count as canonical headedness or as a fusion construction. \textit{CGEL}'s analysis is the latter \citep[Ch.~5 §9]{huddleston2002}. In cases like \mention{I'll take some}, the Det and Head functions of the NP fuse, and a single DP (headed by the determinative \mention{some}) realizes the fused function. Function fusion, not category fusion: \mention{some} remains a determinative; the two syntactic functions it would otherwise fill separately are realized by a single phrase in this construction. The general apparatus is developed by \textcite{Payne2007} as fusion of functions, a first-class theoretical construct with six defining properties.

A first point about head-like behaviour can be separated from the fusion question. In partitives like those in \ref{ex:partitive-headhood}, the determinative is the structural pivot of the matrix NP:
\ea\label{ex:partitive-headhood}
\ea \textit{some of mine}
\ex \textit{some of Kim's}
\ex \textit{many of them}
\ex \textit{all of those}
\z
\z
The \mention{of}-phrase is licensed by the quantificational pivot, exactly as in \textit{some of the books}. \textit{CGEL} itself analyses such partitive cases as fused Det-Head structures rather than as plain canonical headedness \citep[Ch.~5 §9]{huddleston2002}, so this isn't yet evidence that the determinative occupies Head function alone. What it does establish is that the determinative is a head-like pivot in the matrix NP — the lexical anchor on which the partitive is built — independently of whether the construction is fused. That bridge is enough for the rest of the argument: the categorial question (is the determinative a noun?) doesn't require settling the structural question (is the construction fused?) in advance.

Headhood doesn't by itself settle the fusion question. A determinative can head the matrix NP and still be analysed under \textit{CGEL} as fused with Det. The two questions should be kept apart: whether determinatives can head NPs at all, and whether the simple one-word cases should be analysed as ordinary headedness or as fusion.

The direct argument against fused Det-Head analysis for the simple cases then comes from a parallel with pronouns. \textit{CGEL} treats pronoun-headed NPs like \textit{Mine is on the table} or \textit{I'll take yours} as having no Det position: the pronoun heads a Nom, the Nom heads the NP, and there's no Det function to fill or fuse. Even genitive pronouns like \mention{mine}, \mention{yours}, \mention{ours} head NPs this way, without a separate Det function — despite their genitive marking, which in other contexts (the \mention{my} set) co-occurs with Det function. Pronouns are nouns that don't require a determiner.

The categorial conclusion the paper draws is independent of the fusion question. Whether \textit{Some left} is analyzed as \mention{some} heading a Nom directly (parallel to \textit{Mine is on the table}) or as a fused Det-Head construction in which a determinative-headed phrase fills both functions, what's doing the heading is a determinative — and \textit{CGEL}'s treatment of \mention{mine}/\mention{Kim's} already shows that fusion, where it applies, is compatible with the head being a noun. The simpler analysis (canonical Head, no fusion) is preferred where it applies, but the central categorial claim that determinatives are nouns doesn't rest on it.

The cases \textit{CGEL} handles under fusion of functions divide for the present analysis. Simple fused-head, where the dependent function is Det and the item filling it is from the determinative class, makes up the bulk of the discussion. Mod-Head fusion, exemplified by AdjPs in \mention{the rich}, \mention{the poor}, \mention{the elderly}, is a smaller residue. Adjectives head AdjPs, and AdjPs canonically fill modifier function within Nom or predicative complement function within clauses. The Mod-Head fusion machinery is what \textit{CGEL} uses to accommodate the small set of cases where an AdjP also fills Head function — a position the AdjP's category doesn't otherwise reach. Items restricted to predeterminer-modifier function (\mention{quite a}, \mention{rather a}, \mention{many a}, \mention{what a}, \mention{such a}) don't enter the comparison at all: they precede indefinite NPs and never head Noms themselves, fused or otherwise. The substantive question is whether the simple cases pattern with Mod-Head fusion or stand apart. (Items like \mention{all}, \mention{both}, \mention{half} fill predeterminer-modifier function in some constructions and are determinatives elsewhere; they fall under the simple-fused-head treatment below.)

Three positive discriminators show that the simple cases pattern like pronoun-headed NPs (where no Det function is at issue) and unlike adjectival Mod-Head fusion (where fusion is well motivated). The discriminators don't refute Det-Head fusion in general — \textit{CGEL} can keep the fusion analysis for partitive and other cases — but they do narrow the construction-types for which fusion is the most parsimonious analysis.

The first discriminator is paradigm productivity. Fused-head use is a broadly productive pattern within the determinative class. The set that admits it includes \mention{some}, \mention{all}, \mention{both}, \mention{many}, \mention{few}, \mention{several}, \mention{most}, \mention{this/these}, \mention{that/those}, \mention{each}, \mention{either}, \mention{neither}, \mention{much}, and \mention{enough}, and the cardinals (\textit{three of the shirts}, \textit{I'll take two}). \textit{CGEL}~\citep[371--2]{huddleston2002} and \textcite[22--23]{Payne2007} identify a small set of lexically defective members: \mention{the} and \mention{a} lack fused-head use entirely, \mention{every} doesn't appear as a fused head, and \mention{no} is restricted to the suppletive \mention{none}. The defective set is small and independently noted; the pattern remains broadly productive across the rest of the class. Adjective Mod-Head fusion is, by contrast, lexically restricted to a small set, mostly adjectives describing classes of people: \mention{the rich}, \mention{the poor}, \mention{the elderly}, \mention{the dead}, \mention{the unemployed}, \mention{the unobtainable}. The pattern is broadly productive within determinatives and narrowly lexicalized within adjectives.

The second discriminator is structural saturation. Simple fused-head NPs require no additional structure. The examples in \ref{ex:fused-saturated} are full NPs filling argument positions with the determinative as the sole element:
\ea\label{ex:fused-saturated}
\ea \textit{Some left.}
\ex \textit{Many came.}
\ex \textit{All agree.}
\z
\z
Mod-Head fusion requires a determiner. \mention{The rich} in \ref{ex:rich-needs-det}~(a) is a full NP, but \mention{rich} in (b), intended as a class-referring NP, isn't:
\ea\label{ex:rich-needs-det}
\ea \textit{The rich get richer.}
\ex \ungram{\textit{Rich get richer.}}
\z
\z
The Mod-Head AdjP can't saturate the NP on its own; the simple fused-head determinative can.

A third discriminator cuts against covert-head alternatives. In \ref{ex:internal-mod}, an AdjP modifier attaches to the determinative as the lexical pivot of the nominal:
\ea\label{ex:internal-mod}
\ea \textit{The lucky few left.}
\ex \textit{The chosen few will survive.}
\ex \textit{The happy few agree.}
\z
\z
The structure has \mention{few} as the head of the Nom and \mention{lucky/chosen/happy} as an AdjP in Mod function, parallel to \textit{the lucky survivors} where \mention{survivors} is the head and \mention{lucky} is Mod. This isn't discriminating between determinatives and Mod-Head adjectives — the Mod-Head residue also takes AdjP modification (\textit{the idle rich}) — but it does cut against two alternative analyses of the determinative itself. Consider a null-head analysis: the determinative is in Det function of an NP whose head is a silent N, with the AdjP attaching to the silent N. Or a D-with-silent-NP analysis: the determinative is a functional head whose complement is a silent NP, with the AdjP inside that NP. Both analyses face the same structural problem. The AdjP occupies the linear position that attributive modifiers occupy relative to lexical heads (\textit{the lucky winners, the lucky few}): immediately after Det, immediately before the lexical pivot. A silent head would need to be positioned after the AdjP, with the determinative in Det function before it. But the determinative itself occupies the post-modifier position (*\textit{the lucky \_\_\_ few}* is not the structure; the structure is \textit{the lucky few}, with the determinative in the lexical-head slot). The AdjP's attachment target is the determinative, not a silent element, because there's no syntactic position left for a silent head given where the determinative appears.

Three discriminators, three alternative analyses constrained. Productivity separates the simple cases from Mod-Head fusion: the two patterns have different distributional profiles across the class. Saturation rules out a covert-determiner analysis: no such determiner is needed for the NP to be complete. Internal modification rules out null-head and silent-NP analyses: the determinative itself is the pivot that modifiers attach to. None of the three diagnostics defeats Det-Head fusion within an NP-headed analysis (\textit{CGEL}'s live alternative); collectively they make the simpler analysis (canonical Head, no fusion) more parsimonious for the simple cases while leaving fusion as a defensible option for partitive and similar cases. The Mod-Head residue retains fusion on the narrow empirical profile set out above.

The reanalysis runs as follows. The pair in \ref{ex:fused-pair}, traditionally treated as having different structures, receives parallel analyses:
\ea\label{ex:fused-pair}
\ea \textit{I'll take some apples.}
\ex \textit{I'll take some.}
\z
\z
\textit{CGEL} analyses (a) as having \mention{some} in Det function of an NP headed by \mention{apples}, and (b) as having \mention{some} in fused Det+Head function of an NP whose head is anaphorically construed. The present analysis preserves the first; the second has \mention{some} as the head N of an NP, exactly parallel to (\ref{ex:dogs-parallel}):
\ea\label{ex:dogs-parallel}
\textit{I'll take dogs.}
\z
The two sentences in \ref{ex:fused-pair}~(b) and \ref{ex:dogs-parallel} have the same structure in object position; the only difference is which sub-category of N heads the NP.

Under this reanalysis, determinatives canonically head NPs in subject, object, and PrepC functions. The Head function is part of the determinative sub-category's profile rather than a function reached via fusion. The membership criterion is satisfied. The supercategory \term{noun} admits all four sub-categories. Table~\ref{tab:diagnostic} summarizes the result.

\begin{table}[h]
\centering
\caption{The membership criterion across the four sub-categories of \term{noun}, with the Mod-Head residue for contrast. The criterion includes both the distributional condition (heads NP in the canonical nominal functions) and the profile condition (Head function is canonical, not reached via fusion). The determinative row is canonical under the reanalysis in §\ref{subsec:fused-head}.}\label{tab:diagnostic}
\begin{tabular}{lllll}
\hline
Sub-category & Subject & Object & PrepC & Head function \\
\hline
Common N & \textit{People left.} & \textit{I see people.} & \textit{with people} & canonical \\
Proper N & \textit{Kim left.} & \textit{I see Kim.} & \textit{with Kim} & canonical \\
Pronoun & \textit{She left.} & \textit{I see her.} & \textit{with her} & canonical \\
Determinative & \textit{Some left.} & \textit{I see some.} & \textit{with some} & canonical \\
\hline
\multicolumn{5}{l}{\emph{Residue} (does not satisfy the criterion):} \\
Adjective (Mod-Head) & \textit{The rich left.} & \textit{I see the rich.} & \textit{with the rich} & via fusion \\
\hline
\end{tabular}
\end{table}

The reanalysis doesn't eliminate fusion altogether. Mod-Head fusion remains for the small lexicalized set in \mention{the rich}, \mention{the poor}, and so on. The fusion-of-functions apparatus is preserved for those cases, where it earns its keep. What the reanalysis does eliminate is the apparatus's application to the simple cases, which constitute most of \textit{CGEL}'s discussion of fused heads. The simplification is substantial but not absolute.

\subsection{Simplification at the determiner slot}\label{subsec:det-uniform}

Once §\ref{subsec:fused-head} is granted and determinatives are nouns, a downstream simplification follows at the determiner slot. \textit{CGEL}'s existing treatment of the determiner function admits several phrase types: determinative phrases (DPs, in Ch.~5 terminology), genitive NPs, and PPs like \mention{about twenty books} \citep[Ch.~5 §4-5]{huddleston2002}. The core disjunction at the noun-headed end is between DP and genitive NP: two primary lexical categories projecting to different phrase types, both functioning as Det. The disjunction is accommodated but not predicted.

Under the present proposal, the two noun-headed phrase types collapse into one. What was analysed as a DP is an NP headed by a determinative-noun; a genitive NP in Det function is an NP headed by a common or proper noun (or, in the dependent genitive cases, a pronoun). The noun-headed phrase functioning as Det is uniformly an NP. PPs still function as Det alongside, but the primary DP/NP disjunction is simplified.

This isn't a notational change alone. Three observations follow.

First, both phrase types contribute definite reference in parallel ways. \textit{The book} and \textit{Kim's book} are both definite NPs; neither admits a further article. The shared definiteness behaviour is unsurprising if both phrases are NPs whose head can carry a definiteness profile.

Second, the one-determiner-per-NP restriction applies uniformly. None of \textit{*the the book}, \textit{*the Kim's book}, \textit{*Kim's the book}, or \textit{*these Kim's books} works. Whatever fills Det function precludes a further determiner — not because articles specifically are blocked, but because the NP's single Det position is already filled.

Third, fused-head licensing extends naturally. Both phrase types can serve as fused heads in the appropriate constructions, as in \ref{ex:fused-licensing}:
\ea\label{ex:fused-licensing}
\ea \textit{I'll take some.}
\ex \textit{I'll take Kim's.}
\z
\z
Under the present proposal both have NPs in object position, headed in one case by a determinative-noun and in the other by a proper noun marked for genitive case. The genitive NP doesn't require a separate fused-head treatment beyond what the determinative case already needs.

The genitive subsystem shouldn't bear more weight than that. Forms like \mention{mine} and \mention{Kim's} matter because they show that fusion, if present, is compatible with nounhood and pronominality. The independent headhood argument for determinatives comes from the partitives in §\ref{subsec:fused-head}, not from the special morphology of the independent genitives.

The simplification is at the phrase-type level, not the feature level. The head N of the outer NP still selects its Det-functioning NP by features (definite versus indefinite, count versus non-count, number agreement); what changes is that the primary noun-headed phrase is uniformly an NP, rather than a DP-or-NP disjunction. The feature-selection work remains but applies to a single phrase type for the noun-headed cases.

The gain is concrete in one place in particular. Under \textit{CGEL}'s disjunctive analysis, the head noun's selection of Det has to be stated as \enquote{accepts DP or NP (genitive)}, with the two alternatives treated as categorially distinct but functionally parallel. The parallelism itself — the fact that both phrase types contribute the same kinds of content (definiteness, countability selection, fused-head licensing) — is accommodated as a coincidence between two different phrase types. Under the proposal, the parallelism isn't coincidental: it's a consequence of both phrases being NPs whose head nouns contribute the relevant features. \textit{CGEL} has to stipulate the parallelism between DP and NP patterns in Det function; the proposal derives it from a single phrase type. This is the kind of case where the previous account has a descriptive stipulation the new one doesn't need.

The simplification doesn't independently establish determinatives as nouns; that's the work of §\ref{subsec:fused-head}. What §\ref{subsec:det-uniform} cashes out is a consequence of having established it. The noun-headed Det disjunction collapses to a single phrase type, and the patterns that the disjunction left unexplained — definiteness contribution, the one-Det-per-NP restriction, fused-head licensing — receive a uniform structural account for the noun-headed cases.

\section{The articles}\label{sec:articles}

The articles \mention{the} and \mention{a} are the worst case for the proposal. On the criterion in §\ref{sec:argument}, a sub-category of \term{noun} qualifies if its items canonically head NPs; \mention{the} and \mention{a} never head NPs in any structure, not even the fused-head use that §\ref{subsec:fused-head} reanalyses for other determinatives. \textit{CGEL} \citep[371--2]{huddleston2002} and \textcite[22--23]{Payne2007} mark the defectiveness explicitly. \textcite{lyons1999} documents its typological character: articles are phonologically weak, semantically minimal, and cross-linguistically non-universal.

The membership criterion is at the sub-category profile level, as stated in §\ref{sec:argument}. The question isn't whether \mention{the} and \mention{a} individually satisfy the distributional condition — by themselves they don't — but whether the sub-category they belong to satisfies it. Lexical defectiveness within a sub-category is familiar in the grammar: the auxiliary verb \mention{must} lacks a past-tense form (the suppletive \mention{had to} fills the gap), the determinative sub-category includes suppletion (\mention{none} for \mention{no} in fused Det-Head use; \textcite[22--23]{Payne2007}), and common nouns show mass-versus-count defectiveness. Three synchronic considerations bear on the articles specifically.

First, \mention{the} and \mention{a} stand in paradigmatic opposition with the demonstratives (\mention{this/these, that/those}) and the quantifiers (\mention{some, any, every, no, each}) within the determinative sub-category. The contrast is systematic in feature-space: \mention{the} is the definiteness-marked member of a set whose other members carry richer featural content (deixis for demonstratives; quantification for quantifiers); \mention{a} is the indefinite member of a set that includes the numeral \mention{one} and the quantifiers. The articles aren't isolated lexemes outside the sub-category; they're paradigmatic participants in it.

Second, \mention{the} and \mention{a} carry exactly the features the determinative sub-category's profile predicts: definiteness marking, count/non-count selection of the head noun, number specification (\mention{a} for singular count). Common nouns don't carry these features on the lexeme; they allow them to be contributed by a determiner. \mention{The} and \mention{a} contribute what the rest of the determinative sub-category contributes to Det function, in the most lexically reduced form.

Third, the articles are nearly confined to Det function. \mention{The} appears outside Det only in marginal constructions like the correlative comparative (\textit{the bigger the better}), where its analysis is contested. The other determinatives occupy Det as their default but alternate productively into Mod and Head. The articles are the sub-category's defective end for non-Det functions: the end-point of a distributional cline the rest of the sub-category still spans. Defectiveness is relative to the sub-category profile, not relative to common-noun profiles.

\mention{Every} warrants separate comment, because it lacks standalone Head use (\textit{*Every left}, \textit{*I invited every}) and because \textcite{spinillo2004reconceptualising} groups it with \mention{the} and \mention{a} in a residual article class rather than with the quantifiers. The paper places \mention{every} with the other quantifiers (\mention{some, any, no, each, all, most}) on semantic and distributional grounds: it contributes universal quantification, selects singular count heads in the same way \mention{each} does, and admits similar AdvP modifiers (\mention{almost every, nearly every}). Its lexical defectiveness for fused-Head use doesn't pull it out of the quantifier sub-group any more than the defectiveness of \mention{the}/\mention{a} pulls them out of the determinatives as a whole. \mention{Every} joins \mention{the, a, no} as a small set of determinatives with distributional attrition at the fused-head end; their shared defectiveness doesn't define a separate sub-class.

The grammaticalization story explains why \mention{the} and \mention{a} are defective in exactly this way. \mention{The} descends from the Old English distal demonstrative \textit{þæt/se/sēo}; \mention{a/an} descends from the Old English numeral \textit{ān} \enquote{one}. \textcite{heinekuteva2002lexicon} document the cross-linguistic paths: DEMONSTRATIVE~>~DEFINITE ARTICLE and NUMERAL \enquote{ONE}~>~INDEFINITE ARTICLE are among the most robust attested trajectories. \textcite{diessel1999} traces the demonstrative-to-article development in detail; \textcite[Ch.~9]{lyons1999} surveys the cross-linguistic picture. The diachronic account explains the pattern of attrition (phonological reduction, loss of head-of-NP distribution, semantic bleaching) without doing the synchronic membership work. The membership argument rests on paradigmatic, featural, and functional contrast synchronically available in the grammar.

The high token frequency of \mention{the} and \mention{a} follows from their functional role, not from prototype centrality. They're the members of the determinative sub-category with the most discourse work to do (definiteness and indefiniteness marking), so they appear in nearly every NP that takes a determiner. Frequency isn't evidence that they define the sub-category; if anything, it's evidence of how far grammaticalization has reduced them. The sub-category's categorial profile is supplied by the demonstratives, quantifiers, and the rest of the membership; the articles are the grammaticalized endpoint of that profile.

The scope of the claim is synchronic English. Languages vary in how far the grammaticalization cline has run and what form the endpoint takes: Scandinavian has suffixal articles (\textit{hus-et} \enquote{house-the}); Bulgarian and Romanian have similar postposed definiteness markers; many languages lack articles entirely. Whether the fourfold-noun proposal generalizes to systems with bound articles, no articles, or richer agreement morphology is a question for separate work.

\section{Nearby proposals}\label{sec:nearby}

The nearby proposals aren't all equally close. The closest CGEL-internal rival keeps the N-headed NP analysis and the category/function distinction but retains determinative as a separate primary category. The closest English dissolve-Det rival is Spinillo's positive repartitioning. The remaining proposals matter as background because they unify the same terrain in different directions or under different architectures. Table~\ref{tab:proposals} compares them along three dimensions: the unification direction, the status of common nouns, and whether English data alone motivate the proposal.

\begin{table}[h]
\centering
\caption{Nearby proposals on the determinative-pronoun-noun overlap.}\label{tab:proposals}
\small
\begin{tabular}{>{\raggedright\arraybackslash}p{3.1cm}>{\raggedright\arraybackslash}p{6.2cm}>{\raggedright\arraybackslash}p{2cm}>{\raggedright\arraybackslash}p{1.5cm}}
\hline
Proposal & Unification move & Common N status & English data only? \\
\hline
Abney 1987 & D heads DP (functional) & complement of D & no \\
\textit{CGEL}; Pullum \& Miller 2022 & NP-headed; determinative retained as separate primary category & coordinate & yes \\
Postal 1966 & pronouns from article + [+Pro] N & source of pronouns & yes \\
Hudson 2004 & determiners $\subset$ pronouns & separate & yes \\
Anderson 1997 & names, pronouns, and determinatives grouped as \{N\} & \{N;P\} separate & no \\
Van Eynde 2003 & Det split across A and N & unaffected & no \\
Spinillo 2000, 2004 & Det dissolved; residual article class (\mention{the}, \mention{a}, \mention{every}) & unaffected & yes \\
Déchaine \& Wiltschko 2002 & pronouns decompose into pro-DP, pro-$\phi$P, pro-NP & unaffected & no \\
\textbf{present paper} & \textbf{four sub-categories of Noun} & \textbf{coordinate} & \textbf{yes} \\
\hline
\end{tabular}
\normalsize
\end{table}

The nearest rival is the NP-headed tradition that stops one step earlier. \textit{CGEL} and \textcite{pullummiller2022nps} keep the N-headed NP analysis and the category/function distinction while retaining determinative as a separate primary category; \textcite{bruening2020nominal} reaches the same NP-headed conclusion in a generative setting. The disagreement is therefore narrow. Once partitives establish determinative headhood and the simple one-word cases are treated as ordinary NP headedness, the remaining question is taxonomic: why should determinatives still be denied nounhood if the projection facts are already shared?

The nearest English dissolve-Det rival is Spinillo. \textcite{spinillo2000determiners,spinillo2004reconceptualising} offer a positive repartitioning in which demonstratives and genitives join pronouns, quantifiers split across classes, and \mention{the}, \mention{a}, and \mention{every} form a residual article class. The present paper agrees with the negative diagnosis and rejects the repartition. Demonstratives stay with determinatives rather than moving to pronouns, and \mention{every} stays with the quantifiers rather than joining the articles. The articles are handled as defective determinatives, not as the residue left over after repartition.

\textcite{vaneynde2003determiner} and \textcite{herslund2008articles} address different evidence. Van Eynde splits determinatives between adjective and noun on agreement and inflectional evidence from Italian and Dutch. Herslund treats definite articles specifically as pronominal heads. Both are useful comparators, but neither supplies the target claim here: a synchronic English argument that keeps the determinative class together. English lacks the rich prenominal agreement morphology that drives Van Eynde's split, and the article-specific move doesn't generalize to the demonstratives and quantifiers.

\textcite{postal1966}, \textcite{hudson2004determiners}, and \textcite{anderson1997notional} unify the same territory in the opposite direction or under a broader notional cover term. Their shared point is that the determiner-pronoun boundary isn't a clean categorial divide. The present paper agrees with that much. It diverges by refusing to reduce one noun-like sub-category to another or to leave common nouns outside the same supercategory. Pronoun and determinative remain distinct sub-clusters, and common nouns remain coordinate with them inside \term{noun}.

\textcite{dechainewiltschko2002} and \textcite{abney1987} move the dispute to a different architecture. Déchaine and Wiltschko decompose pronouns into pro-DP, pro-$\phi$P, and pro-NP. Abney makes D the functional head of the nominal phrase. Those proposals matter because they also reject a simple pronoun/determinative/common-noun partition. They aren't close rivals here because the present paper keeps the category/function distinction and the N-headed NP analysis that both approaches abandon.

\section{Sub-cluster combinatorics}\label{sec:subcluster}

The supercategory claim concerns NP projection and external distribution. Within \term{noun}, the sub-clusters differ in internal morphosyntax, modifier profiles, and discourse behaviour. That is what the sub-category labels are for. The four sub-categories are distinct for the same kinds of reasons that other grammatical sub-clusters are distinct.

Three sub-cluster differences matter for the present paper.

First, determinative-headed NPs accept AdvP modifiers that common-noun-headed NPs don't:
\ea\label{ex:advp-det}
\ea \textit{almost every problem}
\ex \textit{hardly any mistakes}
\ex \textit{very nearly all responses}
\z
\z
Common-noun-headed NPs take AdjP modifiers instead (\textit{a genuine professional}, \textit{an experienced teacher}). The asymmetry is real, but it marks a sub-cluster difference rather than a supercategory boundary. \textcite[13--14]{payne2010} treat the same determinative across two functions without changing its category: \mention{few} is a determinative in \textit{few doctors} (Det function) and in \textit{the few doctors} (Mod function). The item's category stays constant; its functional profile includes both Det and Mod. On the present proposal the same item is also a noun, so the modifier pattern belongs to the determinative sub-category's profile within \term{noun}.

Second, common and proper nouns usually carry more descriptive content on the lexeme. Pronouns and determinatives are more reduced, with reference tracked through discourse mechanisms such as anaphora, deixis, and quantification rather than lexical description. The contrast isn't absolute — \textit{Many are called, few are chosen} is generic, and deictic uses of demonstratives aren't anaphoric — but the tendency is strong enough to place the two closed-class sub-categories at the semantically reduced end of the supercategory.

Third, the one-Determiner-per-NP restriction is functional, not categorial. An NP contains at most one Det function, but the determinative sub-category can contribute more than one lexical item to a single NP in different functions. The standard case is \textit{the many people}, where \mention{the} fills Det and \mention{many} fills Mod. Both are determinatives; both are nouns on the present account. The functional restriction is preserved without relabeling either item. \textcite{reynolds2026numerals} makes the same point for numeral factors: \textit{two hundred books} has \mention{two} in Mod and \mention{hundred} as Head, with the NP taking a single Det (null in the example, overt in \textit{the two hundred books}).

Empirical support for the sub-cluster structure comes from \textcite{reynolds2021}. An energy-distance clustering analysis applied to a 138-word-form by 232-feature matrix of determinatives and pronouns \citep{reynolds2021matrix} returns a 93\% correspondence with \textit{CGEL}'s sub-category boundaries. The result confirms that the two sub-clusters are distributionally distinct at a level a non-parametric test can recover. The present proposal doesn't dispute that distinctness. It locates it at the sub-category level under a common supercategory rather than across primary lexical categories.

Cardinal numerals are the cleanest case of cross-sub-category lexical alternation, and the proposal simplifies the picture \citet{reynolds2026numerals} draws. Reynolds analyses cardinals as spanning the primary categories determinative and noun depending on use: cardinal-as-determinative in \textit{ten men}, cardinal-as-proper-noun in \textit{Room 101}, cardinal-as-common-noun in \textit{tens of pens}. Under the present account, this cross-categorial span becomes a within-supercategory span: cardinals are nouns in all three uses, with the variation falling among the sub-categories of \term{noun}. The cross-category alternation Reynolds documents is recast as a sub-category alternation within \term{noun}. No fifth sub-category is needed, and no cross-primary-category ambiguity has to be stipulated for cardinal lexemes.

Sub-cluster differences are expected in any non-trivial taxonomy. Common nouns and proper nouns differ in article distribution, pluralization, and descriptive content. Pronouns differ from both in case inflection and paradigmatic closure. Determinatives differ in modifier selection, semantic reduction, and restriction to Det and Mod functions. The sub-category labels track those differences. The supercategory label tracks what the four share: the projection of an NP that heads an argument-bearing constituent.

\section{Conclusion}\label{sec:conclusion}

English has four sub-categories of \term{noun}: common noun, proper noun, pronoun, and determinative. \textit{CGEL} admits three; this paper adds the fourth. The categorial claim doesn't depend on abolishing fused Det-Head analysis. \textit{CGEL}'s existing treatment of independent genitives (\mention{mine}, \mention{Kim's}) already shows that fusion, where it applies, is compatible with noun heads; the paper extends this to the determinative class. Two stages of evidence support the extension. Partitives establish determinatives as head-like pivots in matrix NPs, independently of whether the construction is fused. Three positive discriminators — paradigm productivity, structural saturation, and internal modification — show that simple one-word cases pattern with non-fused noun-headed NPs and unlike adjectival Mod-Head fusion. \textit{CGEL} can keep fused Det-Head for partitive and similar cases; the items doing the fusing are nouns either way.

There is precedent for the move within the same framework. \textit{CGEL} already treats auxiliaries as a sub-category of verb \citep{pullumwilson1977} and admits pronouns to \term{noun} despite their distinct morphosyntax. Extending the sub-categorization of \term{noun} to include determinatives applies the same kind of move to a different closed functional class. The proposal preserves the category/function distinction and the N-headed NP analysis; what changes is the supercategory assignment of determinatives.

Two simplifications follow. The noun-headed phrases that function as Det — what \textit{CGEL} treats as DPs (headed by determinatives) and genitive NPs (headed by common, proper, or pronoun) — collapse into a single phrase type once determinatives are nouns. PPs continue to function as Det alongside, but the primary DP/NP disjunction is simplified. The fusion-of-functions apparatus developed in \textcite{Payne2007} narrows: it drops out of the simple one-word cases (where the discriminators support canonical Head) and is retained for the narrower Mod-Head, Predet-Head, and partitive cases where it remains empirically motivated. The articles \mention{the} and \mention{a} are handled as defective members of a sub-category whose profile otherwise satisfies the criterion; their distributional attrition is the endpoint of a grammaticalization cline whose sources (demonstratives and the numeral \mention{one}) are already nominal on the present analysis.

Sub-cluster distinctions are preserved. Common nouns and proper nouns differ in article distribution and descriptive content; pronouns differ from both in case inflection and paradigmatic closure; determinatives differ in modifier selection, semantic reduction, and distributional restriction to Det and Mod functions. What the four sub-categories share — the projection of an NP filling subject, object, and PrepC function — is what the supercategory tracks.

The scope is synchronic English. Cross-linguistic generalization depends on how far the grammaticalization cline has run and what form the endpoint takes in each language. Whether the fourfold-noun architecture extends to systems with bound articles, no articles, or richer nominal agreement morphology is a question for separate work.

\section*{Acknowledgements}

This paper was drafted with the assistance of large language models. All content has been reviewed and revised by the author, who takes full responsibility for the final text.

\newpage
\printbibliography

\end{document}
